\chapter{D.E.R.}\label{cap:resultados}

% \lipsum[73]

\section{Base de Dados}
\begin{figure}[htpb]
   \centering
   \includegraphics[scale=.05]{figs/der}
   \caption{Proposta de DER para o problema.}
   \label{fig:der}
\end{figure}


\section{Dicionário de Dados}

% \begin{flushleft}
% \begin{longtable}{p{2.5cm} p{4.5cm} p{1.2cm} p{5.5cm} p{3.5cm}}
% \caption{Dicionário de Dados} \\
% \hline
% \textbf{Entidade} &
% \textbf{Campo} &
% \textbf{Tipo} &
% \textbf{Descrição} &
% \textbf{Restrições/Observações} \\
% \hline
% \endfirsthead

% \hline
% \textbf{Entidade} &
% \textbf{Campo} &
% \textbf{Tipo} &
% \textbf{Descrição} &
% \textbf{Restrições/Observações} \\
% \hline
% \endhead

% \textbf{ELEICAO} &
% \texttt{cd\_eleicao} &
% INT &
% Identificador da eleição. &
% Identificador da eleição no TSE. \\

% &
% \texttt{ano} &
% INT &
% Ano em que a eleição ocorreu. &
% -- \\

% &
% \texttt{turno} &
% INT &
% Número do turno da eleição. &
% 1 ou 2, conforme a eleição. \\

% &
% \texttt{tipo} &
% INT &
% Código do tipo de eleição. &
% Código utilizado pelo TSE. \\

% &
% \texttt{ds\_eleicao} &
% VARCHAR &
% Descrição ou denominação da eleição. &
% -- \\

% &
% \texttt{quantidade\_de\_vagas} &
% INT &
% Quantidade de vagas disponíveis no pleito para determinado cargo. &
% -- \\

% \textbf{CARGO} &
% \texttt{cod\_cargo} &
% INT &
% Identificador do cargo eletivo. &
% \textbf{PK}. \\

% &
% \texttt{ds\_cargo} &
% VARCHAR &
% Descrição do cargo conforme a base eleitoral. &
% -- \\

% &
% \texttt{nome} &
% VARCHAR &
% Nome padronizado do cargo. &
% -- \\

% \textbf{MUNICIPIO} &
% \texttt{cd\_municipio} &
% INT &
% Código identificador do município na base eleitoral. &
% \textbf{PK}. \\

% &
% \texttt{nm\_municipio} &
% VARCHAR &
% Nome do município. &
% -- \\

% &
% \texttt{cd\_ibge} &
% INT &
% Código do município segundo o IBGE. &
% FK para identificação territorial no IBGE. \\

% \textbf{UF} &
% \texttt{cd\_ibge} &
% INT &
% Código identificador da unidade federativa segundo o IBGE. &
% \textbf{PK}. \\

% &
% \texttt{sigla} &
% VARCHAR &
% Sigla da unidade federativa. &
% Ex.: PI, DF, SP. \\

% \textbf{POLITICO} &
% \texttt{cpf} &
% VARCHAR &
% CPF utilizado para identificar o político/candidato. &
% \textbf{PK}. \\

% &
% \texttt{nome} &
% VARCHAR &
% Nome do político. &
% -- \\

% &
% \texttt{ds\_grau\_de\_instrucao} &
% VARCHAR &
% Grau de instrução declarado pelo candidato. &
% Informação eleitoral. \\

% &
% \texttt{dt\_nascimento} &
% DATE &
% Data de nascimento do político. &
% -- \\

% \textbf{MANDATO} &
% -- &
% -- &
% Entidade destinada a representar o mandato exercido por um político. &
% Os atributos devem ser definidos conforme os dados necessários para representar o período e o cargo do mandato. \\

% \textbf{PARTIDO} &
% \texttt{numero} &
% INT &
% Número eleitoral do partido. &
% \textbf{PK}. \\

% &
% \texttt{sigla} &
% VARCHAR &
% Sigla do partido político. &
% -- \\

% &
% \texttt{nome} &
% VARCHAR &
% Nome do partido político. &
% -- \\

% &
% \texttt{vies\_politico} &
% FLOAT &
% Valor numérico utilizado para representar o viés político atribuído ao partido. &
% Métrica calculada pelo sistema. \\

% \textbf{POLITICO\_ELEICAO} &
% \texttt{sq\_candidato} &
% INT &
% Identificador único da candidatura em uma determinada eleição. &
% \textbf{PK}; identificador proveniente da base do TSE. \\

% \textbf{ZONA\_ELEITORAL} &
% \texttt{numero} &
% INT &
% Número identificador da zona eleitoral. &
% Pode ser combinado com o município para identificação da zona. \\

% \textbf{VOTACAO\_PARTIDO\_MUNZONA} &
% \texttt{qt\_votos\_legenda\_validos} &
% INT &
% Quantidade de votos válidos atribuídos à legenda partidária na zona eleitoral. &
% -- \\

% &
% \texttt{qt\_total\_votos\_leg\_validos} &
% INT &
% Quantidade total de votos válidos de legenda contabilizados. &
% -- \\

% &
% \texttt{qt\_votos\_nominais\_validos} &
% INT &
% Quantidade de votos nominais válidos recebidos pelos candidatos. &
% -- \\

% \textbf{VOTACAO\_CANDIDATO\_MUNZONA} &
% \texttt{qtd\_votos\_nominais} &
% INT &
% Quantidade de votos nominais recebidos pelo candidato em determinada zona eleitoral. &
% -- \\

% \textbf{PERFIL\_ELEITORADO\_MUNZONA} &
% \texttt{ds\_grau\_escolaridade} &
% VARCHAR &
% Grau de escolaridade dos eleitores. &
% Categoria de escolaridade utilizada pelo TSE. \\

% &
% \texttt{qtd\_eleitores} &
% INT &
% Quantidade de eleitores pertencentes à categoria de escolaridade informada. &
% -- \\

% \textbf{VOTACAO\_MUNICIPIO} &
% \texttt{qt\_comparecimento} &
% INT &
% Quantidade de eleitores que compareceram para votar. &
% -- \\

% &
% \texttt{qt\_brancos} &
% INT &
% Quantidade de votos brancos. &
% -- \\

% &
% \texttt{qt\_nulos} &
% INT &
% Quantidade de votos nulos. &
% -- \\

% &
% \texttt{qt\_abstencoes} &
% INT &
% Quantidade de eleitores que não compareceram à votação. &
% -- \\

% \textbf{RECEITA\_DE\_CAMPANHA} &
% \texttt{cd\_receita} &
% INT &
% Identificador do registro de receita de campanha. &
% \textbf{PK}. \\

% &
% \texttt{vr\_receita} &
% FLOAT &
% Valor monetário da receita de campanha. &
% Representado em reais. \\

% &
% \texttt{ds\_receita} &
% VARCHAR &
% Descrição ou classificação da origem da receita. &
% -- \\

% \textbf{DESPESA\_DE\_CAMPANHA} &
% \texttt{cd\_despesa} &
% INT &
% Identificador do registro de despesa de campanha. &
% \textbf{PK}. \\

% &
% \texttt{sq\_despesa} &
% VARCHAR &
% Identificador sequencial do registro da despesa na base eleitoral. &
% -- \\

% &
% \texttt{ds\_despesa} &
% VARCHAR &
% Descrição ou classificação da despesa de campanha. &
% -- \\

% &
% \texttt{vr\_despesa\_contratada} &
% FLOAT &
% Valor contratado referente à despesa de campanha. &
% Representado em reais. \\

% \textbf{BEM} &
% \texttt{vr\_bem\_candidato} &
% FLOAT &
% Valor declarado dos bens pertencentes ao candidato. &
% Representado em reais. \\

% \textbf{FORNECEDOR} &
% \texttt{cpf\_cnpj\_fornecedor} &
% VARCHAR &
% CPF ou CNPJ do fornecedor relacionado à despesa de campanha. &
% Identificador fiscal do fornecedor. \\

% &
% \texttt{numero} &
% INT &
% Identificador do fornecedor no modelo de dados. &
% \textbf{PK}. \\

% &
% \texttt{ds\_cnae\_fornecedor} &
% VARCHAR &
% Código ou descrição da atividade econômica do fornecedor segundo a CNAE. &
% -- \\

% \textbf{MALHA\_MUNICIPIO} &
% \texttt{geometria} &
% JSON &
% Representação geométrica dos limites territoriais do município. &
% Pode armazenar geometria em formato compatível com dados geoespaciais, como GeoJSON. \\

% \textbf{VIES\_MUNICIPIO} &
% \texttt{ano\_eleitoral} &
% INT &
% Ano da eleição utilizada para o cálculo do viés político municipal. &
% -- \\

% &
% \texttt{vies\_calculado} &
% FLOAT &
% Valor calculado que representa o viés político observado no município. &
% Métrica calculada a partir dos resultados eleitorais. \\

% \textbf{VOTACAO\_PARTIDO\_MUNICIPIO} &
% \texttt{qt\_votos\_validos} &
% INT &
% Quantidade de votos válidos recebidos pelo partido no município. &
% -- \\

% &
% \texttt{ano\_eleicao} &
% INT &
% Ano da eleição à qual a votação pertence. &
% -- \\

% \textbf{IBGE\_CENSO\_MUNICIPIO} &
% \texttt{idade\_media\_populacao} &
% FLOAT &
% Idade média da população do município. &
% Dado proveniente do Censo/IBGE. \\

% &
% \texttt{ano} &
% INT &
% Ano de referência do dado censitário. &
% -- \\

% \textbf{INDICES\_MUNICIPIO} &
% \texttt{ano} &
% INT &
% Ano de referência dos indicadores municipais. &
% -- \\

% &
% \texttt{idhm} &
% FLOAT &
% Índice de Desenvolvimento Humano Municipal. &
% Indicador socioeconômico. \\

% &
% \texttt{pib} &
% FLOAT &
% Produto Interno Bruto do município. &
% Valor monetário, conforme a fonte utilizada. \\

% &
% \texttt{mediana\_idade} &
% INT &
% Mediana da idade da população municipal. &
% -- \\

% &
% \texttt{indice\_envelhecimento} &
% INT &
% Indicador que representa a relação entre a população idosa e a população jovem. &
% Definição deve seguir a metodologia da fonte. \\

% \hline
% \end{longtable}
% \end{flushleft}
\begingroup
\small
\setlength{\tabcolsep}{3pt}
\begin{longtable}{>{\raggedright\arraybackslash}p{4.3cm} >{\raggedright\arraybackslash}p{7.5cm} >{\raggedright\arraybackslash}p{2.8cm}}
\caption{Dicionário de Dados -- Estrutura}
\label{tab:dicionario_dados} \\

\hline
\textbf{Entidade} &
\textbf{Campo} &
\textbf{Tipo} \\
\hline
\endfirsthead

\hline
\textbf{Entidade} &
\textbf{Campo} &
\textbf{Tipo} \\
\hline
\endhead

\textbf{ELEICAO} &
\texttt{cd\_eleicao} &
INT \\

&
\texttt{ano} &
INT \\

&
\texttt{turno} &
INT \\

&
\texttt{tipo} &
INT \\

&
\texttt{ds\_eleicao} &
VARCHAR \\

&
\texttt{\seqsplit{quantidade\_de\_vagas}} &
INT \\

\textbf{CARGO} &
\texttt{cod\_cargo} &
INT \\

&
\texttt{ds\_cargo} &
VARCHAR \\

&
\texttt{nome} &
VARCHAR \\

\textbf{MUNICIPIO} &
\texttt{cd\_municipio} &
INT \\

&
\texttt{nm\_municipio} &
VARCHAR \\

&
\texttt{cd\_ibge} &
INT \\

\textbf{UF} &
\texttt{cd\_ibge} &
INT \\

&
\texttt{sigla} &
VARCHAR \\

\textbf{POLITICO} &
\texttt{cpf} &
VARCHAR \\

&
\texttt{nome} &
VARCHAR \\

&
\texttt{\seqsplit{ds\_grau\_de\_instrucao}} &
VARCHAR \\

&
\texttt{dt\_nascimento} &
DATE \\

\textbf{MANDATO} &
-- &
-- \\

\textbf{PARTIDO} &
\texttt{numero} &
INT \\

&
\texttt{sigla} &
VARCHAR \\

&
\texttt{nome} &
VARCHAR \\

&
\texttt{vies\_politico} &
FLOAT \\

\textbf{\seqsplit{POLITICO\_ELEICAO}} &
\texttt{sq\_candidato} &
INT \\

\textbf{\seqsplit{ZONA\_ELEITORAL}} &
\texttt{numero} &
INT \\

\textbf{\seqsplit{VOTACAO\_PARTIDO\_MUNZONA}} &
\texttt{\seqsplit{qt\_votos\_legenda\_validos}} &
INT \\

&
\texttt{\seqsplit{qt\_total\_votos\_leg\_validos}} &
INT \\

&
\texttt{\seqsplit{qt\_votos\_nominais\_validos}} &
INT \\

\textbf{\seqsplit{VOTACAO\_CANDIDATO\_MUNZONA}} &
\texttt{\seqsplit{qtd\_votos\_nominais}} &
INT \\

\textbf{\seqsplit{PERFIL\_ELEITORADO\_MUNZONA}} &
\texttt{\seqsplit{ds\_grau\_escolaridade}} &
VARCHAR \\

&
\texttt{qtd\_eleitores} &
INT \\

\textbf{\seqsplit{VOTACAO\_MUNICIPIO}} &
\texttt{\seqsplit{qt\_comparecimento}} &
INT \\

&
\texttt{qt\_brancos} &
INT \\

&
\texttt{qt\_nulos} &
INT \\

&
\texttt{qt\_abstencoes} &
INT \\

\textbf{\seqsplit{RECEITA\_DE\_CAMPANHA}} &
\texttt{cd\_receita} &
INT \\

&
\texttt{vr\_receita} &
FLOAT \\

&
\texttt{ds\_receita} &
VARCHAR \\

\textbf{\seqsplit{DESPESA\_DE\_CAMPANHA}} &
\texttt{cd\_despesa} &
INT \\

&
\texttt{sq\_despesa} &
VARCHAR \\

&
\texttt{ds\_despesa} &
VARCHAR \\

&
\texttt{\seqsplit{vr\_despesa\_contratada}} &
FLOAT \\

\textbf{BEM} &
\texttt{\seqsplit{vr\_bem\_candidato}} &
FLOAT \\

\textbf{FORNECEDOR} &
\texttt{\seqsplit{cpf\_cnpj\_fornecedor}} &
VARCHAR \\

&
\texttt{numero} &
INT \\

&
\texttt{\seqsplit{ds\_cnae\_fornecedor}} &
VARCHAR \\

\textbf{\seqsplit{MALHA\_MUNICIPIO}} &
\texttt{geometria} &
JSON \\

\textbf{\seqsplit{VIES\_MUNICIPIO}} &
\texttt{ano\_eleitoral} &
INT \\

&
\texttt{vies\_calculado} &
FLOAT \\

\textbf{\seqsplit{VOTACAO\_PARTIDO\_MUNICIPIO}} &
\texttt{\seqsplit{qt\_votos\_validos}} &
INT \\

&
\texttt{ano\_eleicao} &
INT \\

\textbf{\seqsplit{IBGE\_CENSO\_MUNICIPIO}} &
\texttt{\seqsplit{idade\_media\_populacao}} &
FLOAT \\

&
\texttt{ano} &
INT \\

\textbf{\seqsplit{INDICES\_MUNICIPIO}} &
\texttt{ano} &
INT \\

&
\texttt{idhm} &
FLOAT \\

&
\texttt{pib} &
FLOAT \\

&
\texttt{mediana\_idade} &
INT \\

&
\texttt{\seqsplit{indice\_envelhecimento}} &
INT \\

\hline
\end{longtable}

\begin{longtable}{>{\raggedright\arraybackslash}p{3.5cm} >{\raggedright\arraybackslash}p{4cm} >{\raggedright\arraybackslash}p{4.3cm} >{\raggedright\arraybackslash}p{3.2cm}}
\caption{Dicionário de Dados -- Descrições e Restrições}
\label{tab:dicionario_dados_descricao} \\

\hline
\textbf{Entidade} &
\textbf{Campo} &
\textbf{Descrição} &
\textbf{Restrições/\newline Observações} \\
\hline
\endfirsthead

\hline
\textbf{Entidade} &
\textbf{Campo} &
\textbf{Descrição} &
\textbf{Restrições/\newline Observações} \\
\hline
\endhead

\textbf{ELEICAO} & \texttt{cd\_eleicao} & Identificador da eleição. & Identificador da eleição no TSE. \\
& \texttt{ano} & Ano em que a eleição ocorreu. & -- \\
& \texttt{turno} & Número do turno da eleição. & 1 ou 2, conforme a eleição. \\
& \texttt{tipo} & Código do tipo de eleição. & Código utilizado pelo TSE. \\
& \texttt{ds\_eleicao} & Descrição ou denominação da eleição. & -- \\
& \texttt{\seqsplit{quantidade\_de\_vagas}} & Quantidade de vagas disponíveis no pleito para determinado cargo. & -- \\

\textbf{CARGO} & \texttt{cod\_cargo} & Identificador do cargo eletivo. & \textbf{PK}. \\
& \texttt{ds\_cargo} & Descrição do cargo conforme a base eleitoral. & -- \\
& \texttt{nome} & Nome padronizado do cargo. & -- \\

\textbf{MUNICIPIO} & \texttt{cd\_municipio} & Código identificador do município na base eleitoral. & \textbf{PK}. \\
& \texttt{nm\_municipio} & Nome do município. & -- \\
& \texttt{cd\_ibge} & Código do município segundo o IBGE. & FK para identificação territorial no IBGE. \\

\textbf{UF} & \texttt{cd\_ibge} & Código identificador da unidade federativa segundo o IBGE. & \textbf{PK}. \\
& \texttt{sigla} & Sigla da unidade federativa. & Ex.: PI, DF, SP. \\

\textbf{POLITICO} & \texttt{cpf} & CPF utilizado para identificar o político/candidato. & \textbf{PK}. \\
& \texttt{nome} & Nome do político. & -- \\
& \texttt{\seqsplit{ds\_grau\_de\_instrucao}} & Grau de instrução declarado pelo candidato. & Informação eleitoral. \\
& \texttt{dt\_nascimento} & Data de nascimento do político. & -- \\

\textbf{MANDATO} & -- & Entidade destinada a representar o mandato exercido por um político. & Os atributos devem ser definidos conforme os dados necessários para representar o período e o cargo do mandato. \\

\textbf{PARTIDO} & \texttt{numero} & Número eleitoral do partido. & \textbf{PK}. \\
& \texttt{sigla} & Sigla do partido político. & -- \\
& \texttt{nome} & Nome do partido político. & -- \\
& \texttt{vies\_politico} & Valor numérico utilizado para representar o viés político atribuído ao partido. & Métrica calculada pelo sistema. \\

\textbf{\seqsplit{POLITICO\_ELEICAO}} & \texttt{sq\_candidato} & Identificador único da candidatura em uma determinada eleição. & \textbf{PK}; identificador proveniente da base do TSE. \\

\textbf{\seqsplit{ZONA\_ELEITORAL}} & \texttt{numero} & Número identificador da zona eleitoral. & Pode ser combinado com o município para identificação da zona. \\

\textbf{\seqsplit{VOTACAO\_PARTIDO\_MUNZONA}} & \texttt{\seqsplit{qt\_votos\_legenda\_validos}} & Quantidade de votos válidos atribuídos à legenda partidária na zona eleitoral. & -- \\
& \texttt{\seqsplit{qt\_total\_votos\_leg\_validos}} & Quantidade total de votos válidos de legenda contabilizados. & -- \\
& \texttt{\seqsplit{qt\_votos\_nominais\_validos}} & Quantidade de votos nominais válidos recebidos pelos candidatos. & -- \\

\textbf{\seqsplit{VOTACAO\_CANDIDATO\_MUNZONA}} & \texttt{\seqsplit{qtd\_votos\_nominais}} & Quantidade de votos nominais recebidos pelo candidato em determinada zona eleitoral. & -- \\

\textbf{\seqsplit{PERFIL\_ELEITORADO\_MUNZONA}} & \texttt{\seqsplit{ds\_grau\_escolaridade}} & Grau de escolaridade dos eleitores. & Categoria de escolaridade utilizada pelo TSE. \\
& \texttt{qtd\_eleitores} & Quantidade de eleitores pertencentes à categoria de escolaridade informada. & -- \\

\textbf{\seqsplit{VOTACAO\_MUNICIPIO}} & \texttt{\seqsplit{qt\_comparecimento}} & Quantidade de eleitores que compareceram para votar. & -- \\
& \texttt{qt\_brancos} & Quantidade de votos brancos. & -- \\
& \texttt{qt\_nulos} & Quantidade de votos nulos. & -- \\
& \texttt{qt\_abstencoes} & Quantidade de eleitores que não compareceram à votação. & -- \\

\textbf{\seqsplit{RECEITA\_DE\_CAMPANHA}} & \texttt{cd\_receita} & Identificador do registro de receita de campanha. & \textbf{PK}. \\
& \texttt{vr\_receita} & Valor monetário da receita de campanha. & Representado em reais. \\
& \texttt{ds\_receita} & Descrição ou classificação da origem da receita. & -- \\

\textbf{\seqsplit{DESPESA\_DE\_CAMPANHA}} & \texttt{cd\_despesa} & Identificador do registro de despesa de campanha. & \textbf{PK}. \\
& \texttt{sq\_despesa} & Identificador sequencial do registro da despesa na base eleitoral. & -- \\
& \texttt{ds\_despesa} & Descrição ou classificação da despesa de campanha. & -- \\
& \texttt{\seqsplit{vr\_despesa\_contratada}} & Valor contratado referente à despesa de campanha. & Representado em reais. \\

\textbf{BEM} & \texttt{\seqsplit{vr\_bem\_candidato}} & Valor declarado dos bens pertencentes ao candidato. & Representado em reais. \\

\textbf{FORNECEDOR} & \texttt{\seqsplit{cpf\_cnpj\_fornecedor}} & CPF ou CNPJ do fornecedor relacionado à despesa de campanha. & Identificador fiscal do fornecedor. \\
& \texttt{numero} & Identificador do fornecedor no modelo de dados. & \textbf{PK}. \\
& \texttt{\seqsplit{ds\_cnae\_fornecedor}} & Código ou descrição da atividade econômica do fornecedor segundo a CNAE. & -- \\

\textbf{\seqsplit{MALHA\_MUNICIPIO}} & \texttt{geometria} & Representação geométrica dos limites territoriais do município. & Pode armazenar geometria em formato compatível com dados geoespaciais, como GeoJSON. \\

\textbf{\seqsplit{VIES\_MUNICIPIO}} & \texttt{ano\_eleitoral} & Ano da eleição utilizada para o cálculo do viés político municipal. & -- \\
& \texttt{vies\_calculado} & Valor calculado que representa o viés político observado no município. & Métrica calculada a partir dos resultados eleitorais. \\

\textbf{\seqsplit{VOTACAO\_PARTIDO\_MUNICIPIO}} & \texttt{\seqsplit{qt\_votos\_validos}} & Quantidade de votos válidos recebidos pelo partido no município. & -- \\
& \texttt{ano\_eleicao} & Ano da eleição à qual a votação pertence. & -- \\

\textbf{\seqsplit{IBGE\_CENSO\_MUNICIPIO}} & \texttt{\seqsplit{idade\_media\_populacao}} & Idade média da população do município. & Dado proveniente do Censo/IBGE. \\
& \texttt{ano} & Ano de referência do dado censitário. & -- \\

\textbf{\seqsplit{INDICES\_MUNICIPIO}} & \texttt{ano} & Ano de referência dos indicadores municipais. & -- \\
& \texttt{idhm} & Índice de Desenvolvimento Humano Municipal. & Indicador socioeconômico. \\
& \texttt{pib} & Produto Interno Bruto do município. & Valor monetário, conforme a fonte utilizada. \\
& \texttt{mediana\_idade} & Mediana da idade da população municipal. & -- \\
& \texttt{\seqsplit{indice\_envelhecimento}} & Indicador que representa a relação entre a população idosa e a população jovem. & Definição deve seguir a metodologia da fonte. \\

\hline
\end{longtable}
\endgroup



% \lipsum[72]

% \section{Considerações Finais}

% \lipsum[74]
